\documentclass[../thesis.tex]{subfiles}

\begin{document}
	Современные лабораторные комплексы активно эксплуатируют программные сервисы для управления оборудованием,
	сбора, хранения, анализа и визуализации данных. По мере развития, увеличивается как число
	сервисов, так и их сложность. Для эффективного и надежного управления комплексом требуются
	соответствующие процессы и инструменты.

	В лаборатории \textnumero 11 Института ядерной физики им. Г.И. Будкера СО РАН программные
	сервисы разрабатываются и внедряются в течение длительного времени различными разработчиками с
	использованием различных технологий и инструментов. Подобный эволюционный подход позволял оперативно
	решать локальные задачи, однако с ростом количества сервисов он перестал масштабироваться и стал
	источником операционных рисков. Для дальнейшего успешного функционирования требуется новая современная
	модель эксплуатации, разработанная с учетом уникальных условий и задач лаборатории.

	Объектом исследования в данной работе является программная инфраструктура лабораторного
	комплекса \textnumero 11 коллайдера ВЭПП-2000, включающая совокупность программных сервисов,
	обеспечивающих управление установкой, мониторинг ее состояния, а также обработку и визуализацию
	экспериментальных данных.

	Предметом исследования являются методы и средства унифицированного управления жизненным циклом
	гетерогенных программных сервисов лаборатории, включая механизмы их определения, конфигурации, запуска,
	мониторинга и предоставления операторского интерфейса в условиях локальной эксплуатации.

	\textbf{Цель работы} состоит повышении качества процессов разработки, внедрения, сопровождения и эксплуатации
	программных сервисов лабораторного комплекса. Приоритет отдается \textbf{практической} применимости
	предложенных решений для \textbf{существующих} проблем.

	\textbf{Практическая значимость} заключается в модернизации процессов разработки, внедрения, сопровождения
	и эксплуатации программных сервисов лаборатории \textnumero 11 ИЯФ. Учет специфик лаборатории
	позволит интегрировать систему в ограниченных условиях без нарушения существующих рабочих
	процессов. Сделанные наблюдения и разработанные решения могут быть использованы в других лабораториях
	со схожими потребностями.
\end{document}